\chapter{คลังรหัสต้นฉบับและองค์ประกอบสกิลจำลองฟิสิกส์ 3 มิติ และ AIoT}

ภาคผนวกนี้รวบรวมข้อมูลจำเพาะ โครงสร้างเมทาดาทา และโค้ดอัลกอริทึมของสกิลปัญญาประดิษฐ์ (Agent Skills) ประจำเล่มที่ 1 อย่างละเอียด เพื่อเป็นคู่มืออ้างอิงเชิงวิศวกรรมซอฟต์แวร์

\section{สกิล: nextgen-physics-virtual-lab-architect}

\begin{teachingplan}{ข้อมูลจำเพาะและเมทาดาทาของสกิล}
\textbf{ชื่อสกิล:} \texttt{nextgen-physics-virtual-lab-architect}\\
\textbf{วัตถุประสงค์:} ออกแบบและสร้างห้องปฏิบัติการฟิสิกส์ 3 มิติเสมือนจริงตามมาตรฐานโอลิมปิกวิชาการ (IJSO Standard)\\
\textbf{ขอบเขตเทคโนโลยี:} Three.js PBR Shaders, Raycasting, SVG Vector Engine, Euclidean GCD
\end{teachingplan}

\subsection{โค้ดแกนกลาง: เอนจินคำนวณเศษส่วนแท้ Euclidean GCD}

\begin{lstlisting}[language=JavaScript, caption={Euclidean GCD Engine สำหรับรักษาความแม่นยำทางฟิสิกส์}]
function gcd(a, b) {
  a = Math.abs(Math.round(a));
  b = Math.abs(Math.round(b));
  while (b > 0) {
    const t = b;
    b = a % b;
    a = t;
  }
  return a;
}

function simplifyFraction(numerator, denominator) {
  if (denominator === 0) return { num: 0, den: 1 };
  const common = gcd(numerator, denominator);
  return {
    num: numerator / common,
    den: denominator / common,
    str: `${numerator / common}/${denominator / common}`
  };
}
\end{lstlisting}

\subsection{โค้ดแกนกลาง: การคำนวณความต้านทานสมมูลวงจรหกเหลี่ยม}

\begin{lstlisting}[language=JavaScript, caption={Hexagonal Circuit Reduction Algorithm}]
function calculateHexEquivalentResistance(nodeP, nodeQ, resistors) {
  // nodeP, nodeQ: Indices 0 to 5 (Nodes a, b, c, d, e, f)
  const totalR = resistors.reduce((sum, val) => sum + val, 0);
  
  // Calculate direct path resistance (r1)
  let r1 = 0;
  let curr = nodeP;
  while (curr !== nodeQ) {
    r1 += resistors[curr];
    curr = (curr + 1) % 6;
  }
  
  // Calculate indirect path resistance (r2)
  const r2 = totalR - r1;
  
  // Parallel equivalence: Req = (r1 * r2) / (r1 + r2)
  const req = (r1 * r2) / totalR;
  const fraction = simplifyFraction(r1 * r2, totalR);
  
  return { r1, r2, req, fraction };
}
\end{lstlisting}

\section{สกิล: ar-simulation-builder}

\begin{teachingplan}{ข้อมูลจำเพาะและเมทาดาทาของสกิล}
\textbf{ชื่อสกิล:} \texttt{ar-simulation-builder}\\
\textbf{วัตถุประสงค์:} สร้างระบบโต้ตอบ 3 มิติไร้สัมผัสผ่านการตรวจจับท่าทางมือด้วย MediaPipe Hands 60 FPS\\
\textbf{ขอบเขตเทคโนโลยี:} MediaPipe Vision, Web Audio API Acoustic Synthesizer, Zero-GC Vector Math
\end{teachingplan}

\subsection{โค้ดแกนกลาง: ระบบตรวจจับการจีบนิ้วและการสังเคราะห์เสียง}

\begin{lstlisting}[language=JavaScript, caption={Pinch Gesture Detection and Web Audio Feedback}]
class TouchlessHandController {
  constructor(audioCtx) {
    this.audioCtx = audioCtx;
    this.isPinching = false;
    this.threshold = 0.05; // Normalized screen space distance
  }

  detectPinch(landmarks) {
    const thumb = landmarks[4];
    const index = landmarks[8];
    const dist = Math.hypot(
      thumb.x - index.x, 
      thumb.y - index.y, 
      thumb.z - index.z
    );

    if (dist < this.threshold && !this.isPinching) {
      this.isPinching = true;
      this.playHapticSound(880, 0.05); // High-pitch click
      return 'PINCH_START';
    } else if (dist >= this.threshold && this.isPinching) {
      this.isPinching = false;
      this.playHapticSound(440, 0.04); // Release click
      return 'PINCH_END';
    }
    return this.isPinching ? 'PINCHING' : 'IDLE';
  }

  playHapticSound(freq, duration) {
    const osc = this.audioCtx.createOscillator();
    const gain = this.audioCtx.createGain();
    osc.frequency.value = freq;
    gain.gain.setValueAtTime(0.15, this.audioCtx.currentTime);
    gain.gain.exponentialRampToValueAtTime(0.001, this.audioCtx.currentTime + duration);
    osc.connect(gain);
    gain.connect(this.audioCtx.destination);
    osc.start();
    osc.stop(this.audioCtx.currentTime + duration);
  }
}
\end{lstlisting}

\section{สกิล: soil-pht-aiot-architect}

\begin{teachingplan}{ข้อมูลจำเพาะและเมทาดาทาของสกิล}
\textbf{ชื่อสกิล:} \texttt{soil-pht-aiot-architect}\\
\textbf{วัตถุประสงค์:} ควบคุมการรับส่งข้อมูลผ่านฮาร์ดแวร์ RS485 Modbus RTU และฝึกสอนโมเดล Edge AI OLS\\
\textbf{ขอบเขตเทคโนโลยี:} Dart/Flutter Serial Port, Modbus CRC16, Ordinary Least Squares
\end{teachingplan}

\subsection{โค้ดแกนกลาง: การคำนวณผลรวมตรวจสอบ Modbus CRC16}

\begin{lstlisting}[language=JavaScript, caption={Modbus CRC16 Verification Function}]
function calculateCRC16(buffer) {
  let crc = 0xFFFF;
  for (let pos = 0; pos < buffer.length; pos++) {
    crc ^= buffer[pos];
    for (let i = 8; i !== 0; i--) {
      if ((crc & 0x0001) !== 0) {
        crc >>= 1;
        crc ^= 0xA001;
      } else {
        crc >>= 1;
      }
    }
  }
  return crc; // 16-bit CRC Checksum
}
\end{lstlisting}
